\renewcommand{\konferenztag}{\donnerstag}
\newtimeslot{9:00}


% 9:00
%
\abstractEins{Falk Zscheile}{OSM und die Kunst die Welt zu ordnen}{Sprachliche und logische Aspekte 
des Taggingschemas in OpenStreetMap}{%
Anders als bei der klassischen Geodatenerfassung gibt es bei OSM keine hierarchisches gegliedertes System zur Kennzeichnung der erfassten Daten. 
Über die Frage des \glqq richtigen\grqq\  Taggings, wird sehr häufig gestritten. Oft verlaufen die Diskussionen ohne Ergebnis im Sande.

Dabei zeigt sich, dass häufig schon bei der Diskussion eines Problems aneinander vorbeigeredet wird, da der Kontext 
des Diskussionspartners nicht berücksichtigt wird. Dies zeigt sich häufig auch in einer mangelnden Auseinandersetzung mit 
Grundprinzipien der Logik. Zudem bleibt unberücksichtigt, dass sich die Tagfindung aufgrund der \glqq on the ground rule\grqq  an 
der Alltagsbedeutung von Begriffen orientiert, in den anschließenden Diskussionen aber eine feste Bedeutung im Sinne 
wissenschaftlicher Definitionen unterstellt wird.
}

\abstractZwei{Otto Dassau und Sören Gebbert}{GRASS-Funktionalität in QGIS nutzen}{}{%
	Seit 2005 ist GRASS GIS in QGIS über das GRASS Plugin integriert und stellt hunderte 
Analysemethoden über QGIS bereit -- ergänzend mittlerweile auch über das Processing und WPS Plugin. 
	
	Wir starten in diesem Vortrag einen Vergleich der drei Möglichkeiten zur Interaktion von 
GRASS und QGIS. Anhand von Beispielen stellen wir die Vor- und Nachteile dar und versuchen aufzuzeigen, 
welche Variante in welcher Situation die beste Lösung darstellt. Wann ist es sinnvoll das Processing Plugin 
zu verwenden, wann sollte man auf das bewährte GRASS Plugin setzen und welche Stärken und Schwachstellen 
bietet die Integration von GRASS Funktionalität über das WPS Plugin.
	
	Abschließend stellen wir die Aktuelle Entwicklung der Interaktion zwischen GRASS und QGIS vor und schauen in die Zukunft.}


\abstractZehn{Andreas Hocevar}{OpenLayers 3}{Stand, Neues und Zukünfiges}{%
	OpenLayers 3 wird vorgestellt, zahlreiche Beispiele zeigen die Verwendung. Unterschiede zur Vorgängerversion 
	werden aufgezeigt und auch wo und warum OpenLayers 3 anders ist. Aktuelle Entwicklungen, wie etwa das OL3-Cesium 
	Project, welches die dritte Dimension für OpenLayers zuänglich macht, werden präsentiert und ein Blick in die 
	zukünftige Entwicklung gewagt. }
%

\newtimeslot{9:30}
% 9:30
%
% Alles weitere, was im lagen Abstract steht und hier nicht auskommentiert ist, ist relativ logisch. 
% Außerdem ist der Rest des Absatzes auch mit inhaltlichen Fehlern gespickt, 
% die man niemanden lesen lassen möchte.
% Originaler Titel im Pentabarf "OSM auf Räder", nicht "OSM auf Rädern"!
%\abstractEins{Nicole Beringer}{OSM auf Räder}{}{Sind die OSM Karten soweit, um den Wettbewerb mit kommerziellen Karten im Autonavigationsbereich aufzunehmen?}
\abstractEins{Nicole Beringer}{OSM auf Räder}{}{%
	Sind die OSM-Karten soweit, um den Wettbewerb mit kommerziellen Karten im Autonavigationsbereich aufzunehmen?	
	Die Software-Unternehmen aus dem automotiven Bereich verfolgen die OSM-Entwicklung mit großer 
	Aufmerksamkeit und stehen auf der Spitze der Implementierung dieser neuen Technologie.	
}

\abstractZwei{Sören Gebbert}{Zeitreihenanalyse mit GRASS GIS}{GRASS das temporale Open Source GIS}{%
	Durch die Integration der Zeit als neue Dimension in GRASS GIS stehen nun über 45 Werkzeuge 
	zur Zeitreihenalayse bereit. 
	Eine Auswahl an Werkzeugen für die Analyse, Verarbeitung und Visualisierung von Raster- und Vektor-Zeitreihen wird vorgestellt.}

\abstractZehn{Dominik Helle}{Ein Duett -- OpenLayers 3 im Zusammenspiel mit AngularJS}{}{%
	AngularJS ist ein mächtiges Open-Source-Framework mit der clientseitige Webanwendungen nach dem MVC-Prinzip erstellt werden können. 
	OpenLayers ist wohl eine der bekanntesten Javascript-Bibliotheken um Karten im Netz nutzen zu können.
	
	Im Vortrag werden Möglichkeiten und Anregungen des Zusammenspiels dieser beiden 
	Komponenten aufgezeigt. 
	Ein erstes Aufeinandertreffen der beiden Komponenten gab es für den Autor dieses Vortrags im Rahmen eines 
	Projektes in dem eine Kartenanwendung umgesetzt wurde. Besonderes Augenmerk lag bei der Entwicklung auf der 
	schnellen und leichten anpassbar- und Wiederverwendbarkeit der Komponenten.}


\sponsorenbox{107_EB+Elektrobit}{0.43\linewidth}{-1em}{%
	\textbf{Silbersponsor}\\
	\begin{otherlanguage}{english}
	Elektrobit (EB) 
	Automotive is an international supplier of 
	embedded software solutions and services for the automotive 
	industry. With over 25 years of experience in the automotive 
	industry, EB Automotive offers flexible, innovative solutions 
	for connected car infrastructure, human machine interfaces (HMI) 
	technologies, navigation, driver assistance, ECU, and software 
	engineering services.
	\end{otherlanguage}
}
%

\newtimeslot{10:00}
% 10:00
%

\abstractEins{Robert Klemm}{Automatisierte OSM Aufbereitung und Analyse von LKW-Mautstrecken in Deutschland}{}{%
	Das beschriebene Tool erstellt mautpflichtigen Strecken anhand von Mautpunkten der Bundesanstalt für Straßenwesen (BASt). 
	Die Daten der BASt bezüglich der Strecken werden so umgewandelt, 
	dass sie in OSM abgelegt werden können. Dazu wird ein Taggingschema vorgeschlagen.
	
	Das beschriebene Tool und die damit erzeugten Mautdaten erleichtern dem Mappern die Arbeit. Für sie 
	validiert das Tool durch Routinganalyse die OSM-Daten, basierend auf dem Taggingschema mit mautbezogenen 
	Routingparametern, wie Fahrzeug-, Achs-, Gewichtsklasse und Betreiber. Es lässt sich beispielsweise die 
	schnellste und günstigste Strecke, bezogen auf die zu entrichtende Straßenmaut, ermitteln und grafisch 
	anzeigen.}

\abstractZwei{Stefan Ziegler}{Kreisbogen in QGIS}{}{Neben den gängigen Vektorelementen (Punkt, Linie, Polygon) 
	unterstützt QGIS neu Kreisbogengeometrien. Es ist jetzt möglich solche Geometrietypen 
	(ST\_CircularString, ST\_CompoundCurve, ST\_CurvePolygon etc.) aus einer Postgis-Datenbank zu lesen, anzuzeigen, zu editieren und wieder zu speichern.}

\abstractZehn{Christian Mayer und Marc Jansen}{GeoExt}{Ein ExtJS basiertes Rich-Client-WebGIS-Framework 
	in Zeiten von AngularJS und Konsorten}{%
	Der Vortrag stellt die neueste Version von GeoExt vor und zeigt auf, welches Handwerkszeug dem Entwickler 
	hier bereitsgestellt wird. Unterschiede zwischen ExtJS und anderen Bibliotheken werden benannt, dies kann 
	als Diskussionsgrundlage für die Wahl einer Bibliothek dienen. Schwerpunkt ist die Betrachtung der zukünftigen 
	Entwicklung von GeoExt.}
%

\newtimeslot{11:15}
% 11:15
%
\abstractEins{Claas Leiner}{Daten aus OSM extrahieren und in QGIS weiterverarbeiten}{Am Beispiel der Verteilung 
	von Windanlagen in Deutschland}{%
	Der Vortrag erläutert die Fragestellung am Beispiel der Verteilung von Windkraftanlagen in Deutschland. A
	ngefangen vom Import der deutschlandweiten OSM-Daten in eine PostGis-Datenbank, über die Abfrage der 
	WKA-Standorte bis hin zur Präsentation der Verteilung in einer farbigen Flächendichtekarte, die aus 
	einer interpolierten Rasteroberfläche erzeugt worden ist, welche die Anzahl der Windanlagen im Umkreis 
	von 20\,km darstellt. Dabei kamen Abfrage- und Geoverarbeitungswerkzeuge aus QGIS sowie das GRASS-Modul 
	zur Spline-Interpolation (v.surf.rst) zur Anwendung.}

\abstractZwei{Mila Frerichs}{Geospatial Ruby}{}{%
	Der Talk gibt einen Überblick darüber, was mit Ruby im Geo Bereich möglich ist. Viele große 
	erfolgreiche Webprojekte sind mit Ruby und dem dazugehörigen Webframework Rails umgesetzt worden. 
	
	 Der Fokus dieses Vortrags liegt auf drei wichtigen Bibliotheken terraformer.rb, rgeo und SimpleTile.}

\abstractZehn{Jörg Thomsen}{MapServer Pro-Tipps}{Das \emph{Do what I want} Know How für den UMN MapServer}{%
	MapServer ist bekannt für seine Zuverlässigkeit, Performance und Stabilität. Gerne wird 
	aber auch über die Konfiguration mit den Mapfiles gestöhnt. Dabei bietet gerade diese Art der 
	Konfiguration Möglichkeiten, die tief in der Dokumentation versteckt sind und die bei kniffligen 
	Augaben wirklich hilfreich sein können. Die meisten Anwender kennen nur die Basis-Konfigurationen, 
	aber wer weiß schon dass es einen eingebauten OpenLayers-Client gibt, der es ermögicht nur mit einem 
	Mapfile eine interakive Karte zu veröffentlichen, wie man Linien mit Pfeilspitzen versieht, Geometrien 
	glättet oder den Zugriff auf MapServer-WMS direkt im Mapfile auf bestimmte IP-Adressen beschränken kann?
	
	Weitere Themen des Vortrags sind PHP-Fassaden, JOINs mit Datei basierten Datenquellen, ColorRanges, Cluster von Punktobjekten, Geomtransform und IP-Zugriffsbeschränkungen.}
%

\newtimeslot{11:45}
% 11:45
%
\abstractEins{Jochen Topf}{Taginfo und wie es die Welt sah}{}{%
	Das OpenStreetMap-Projekt organisiert seine Daten über ein offenes Taggingschema. 
	Jeder kann neue Objekte und Tags erfinden und in der OSM-Datenbank speichern. 
	Dieser Vortrag stellt den Web-Dienst Taginfo vor, der Informationen zu Tags aus 
	verschiedenen Quellen darstellt und Mappern so hilft, die Übersicht zu behalten. 
	Der Vortrag geht sowohl auf die alltägliche Nutzung als auch auf weniger bekannte Seiten des Dienstes ein.}

\abstractZwei{Alexey Valikov}{Jsonix: OGC Web Services in JSON}{}{%
	JSON hat wahrscheinlich XML schon längst als Lingua franca ersetzt. 
	Es ist viel leichtgewichtiger und einfacher zu verwenden als XML, vor 
	allem in den JavaScript-basierten Web-Apps.
	
	Das Web GIS Umfeld wird von JavaScript-Bibliotheken wie OpenLayers und Leaflet dominiert. 
	Für die gehört JSON sowieso zur Muttersprache. Aber die OGC-Standards sind fast alle XML-basiert 
	und durch XML Schemata spezifiziert.
	
	Jsonix ist eine Open-Source Bibliothek für die XML<->JS Konvertierung, die genau das möglich macht. 
	Mit Jsonix kann man ein XML Schema nehmen und daraus eine Mapping-Datei erzeugen.
	Dieser Vortrag gibt eine Überblick von Jsonix und zeigt es Live in WMS, WFS, CSW sowie OL3 WPS Client Demos vor.}

\abstractZehn{Oliver Tonnhofer}{Mapnik oder MapServer}{Wer macht die schönsten Kartenß}{%
	Mit Mapnik und MapServer stehen zwei OpenSource Kartenrenderer zur Verfügung, die in Geschwindigkeit, 
	Funktionsumfang und Bildqualität kaum Wünsche übrig lassen. Aber welche Software nehme ich für mein Projekt?
	
	Der Vortrag geht auf die kleinen und großen Unterschiede zwischen Mapnik und MapServer ein. Für welche 
	Einsatzzwecke ist MapServer besser geeignet? Was kann Mapnik besonders gut? Wie können die Renderer in 
	Anwendungen und Server integriert werden? Gibt es überhaupt nennenswerte Unterschiede?}
%

\newtimeslot{12:15}
% 12:15
%
\abstractEins{Peter Barth}{Crowd-Sourced Elevation}{DIY DEM}{%	
	Digitale Höhenmodelle können für eine Vielzahl von Anwendungen im OpenStreetMap-Umfeld verwendet werden. 
	In Renderstilen werden sie z.\,B. für Schummerungen oder Höhenlinien verwendet. 
	Aber auch bei Routinganwendungen gibt es verschiedenste Gründe, das Geländemodell mit zu beachten: 
	Fahrradfahrer, Rollstuhlfahrer aber auch Energiesparen bei Elektromobilität.
	
	Mit SRTM steht ein genaues Höhenmodell mit nahezu globaler Abdeckung gemeinfrei zur Verfügung und 
	erst kürzlich wurde auch die Version mit einer Auflösung von einer Bogensekunde freigegeben. Aber reicht das?
	
	Nach einem Überblick über das Thema Höhenmodelle und die verschiedenen Datenquellen wird anhand von Experimenten und 
	deren Auswertung gezeigt, wie man mit einem handelsüblichen Smartphone selbst ein Höhenmodell erstellen 
	kann. Eignen sich die Methoden für Crowdsourcing und wie gut ist das Ergebnis verglichen mit SRTM?}

\abstractZwei{Christian Mayer}{Serverseitiges JavaScript und GIS}{Mit Node.js und npm die Welt beherrschen}{Der 
	Vortrag gibt eine allgemeine Einführung in die Node.js-Welt und zeigt wie mit Node.js geo-relevante Probleme, 
	wie das Einlesen von GIS-Datenformaten sowie die Verarbeitung von Geodaten möglich wird. Ebenso wird eine 
	Übersicht von nützlichen Node.js-Modulen aus dem GIS-Bereich sowie Beispielanwendungen vorgestellt.}

\abstractZehn{Daniel Koch}{GeoServer in action}{Fortgeschrittene Möglichkeiten beim Einsatz des GeoServers}{%
	Der GeoServer ist ein weithin bekannter und mächtiger OpenSource Kartenserver. Sofern man die Umgebung 
	eingerichtet hat, ist sowohl die Installation als auch die Konfiguration von ersten WMS- und WFS-Layern sehr einfach. 
	In diesem Vortrag wird auf die typischen Anforderungen eines \glqq GeoServers in action\grqq\ eingegangen. Der Vortrag fokussiert 
	sich also weniger auf die \glqq Out-of-the-box\grqq{}-Verwendung des GeoServers, sondern beleuchtet vielmehr fortgeschrittene 
	Möglichkeiten beim Einsatz dieser Software.
	
	Themen des Vortrags sind das Kompilieren, die Verwendung der REST-Schnittstelle, Extensions, Performance-Tuning, GeoWebCache-Einsatz, Troubleshooting und Stolperfallen.}
%

\newtimeslot{13:45}
% 13:45
%
\abstractAula{Arnulf Christl}{FOSSGIS-Zukunftswerkstatt}{Menschen begeistern im FOSSGIS-Verein mitzuwirken
	}{%
	Wer steckt hinter dem FOSSGIS e.V.? Um diese Frage zu beantworten veranstaltet der FOSSGIS e.V eine Zukunftswerkstatt. 
	Erfahren Sie mehr über Ziele, Strukturen und Menschen im Verein. 
	Lassen Sie sich begeistern von der Idee gemeinsam etwas zu bewegen. }
\vspace{-1em}
\sponsorenbox{101_52n}{0.38\linewidth}{-1em}{%
	\textbf{Silbersponsor}\\
	52$^\circ$\,North ist ein offenes in\-ternationales Forschungs- und Ent\-wick\-lungsnetzwerk im Bereich der Geoinformatik. 
	Partner aus Forschung, Industrie und Verwaltung  arbeiten aktiv an  Innovationsthemen, die aus den 
	Forschungslabors in die breite praktische Anwendung gebracht werden. Die 52$^\circ$\,North GmbH dient dem 
	Netzwerk als Geschäfts- und Koordinierungsstelle sowie als Service-Center. Kernaufgabe ist die 
	Entwicklung neuer Konzepte und Technologien in der Geoinformatik, insbesondere in den Bereichen 
	Sensor Web und Web basiertes Geoprocessing. Darüber hinaus bietet 52$^\circ$\,North auf Basis seiner 
	Erfahrungen und Technologien Dienstleistungen in den Bereichen Consulting und Software-Entwicklung an.}
\newpage\noindent%\sponsorenbox{105_terrestris}{0.45\linewidth}{-1em}{%
%	terrestris bietet Dienstleistungen und Produkte mit Open Source Software an und 
%	ist dabei insbesondere auf Geodateninfrastrukturen, webbasierte Informationssysteme und WebMapping-Anwendungen fokussiert, 
%	die neben 2D auch 3D Daten enthalten und auch für mobile Endgeräte optimiert sind.
%	
%	Wir orientieren uns dabei an den Anforderungen unserer Kunden und realisieren auf 
%	den jeweiligen Bedarf zugeschnittene Lösungen.
%	Neben der Anwendungsentwicklung können unsere Kunden Support und Schulungen für 
%	verschiedene Open Source Softwaretools, z. B. PostgreSQL/PostGIS, UMN Mapserver, 
%	GeoServer, SHOGun, OpenLayers, GeoExt, ExtJS und QGIS, erhalten.
%	
%	National und international gut vernetzt bieten wir zusammen mit unseren Partnerfirmen
%	Wartung und Support zu Open Source Produkten/Projekten wie SHOGun, GeoServer oder OpenLayers an.
%	Diese Rahmenbedingungen erlauben uns unvoreingenommen Lösungen zu entwickeln und 
%	anzubieten, die tatsächlich alle Anforderungen unserer Kunden zufrieden stellen.
%}
%
\vfill
\setlength{\fboxsep}{4.5pt}
\noindent\fcolorbox{gray}{white}{\parbox{\fboxwidth}{%
	\begin{wrapfigure}{r}{0.3\textwidth}
		\vspace{-1em}
		\centering 
		\includegraphics[width=\linewidth]{105_terrestris}
		\vspace{-1em}
	\end{wrapfigure}	
	\textbf{Silbersponsor}\\	
	\noindent terrestris bietet Dienstleistungen und Produkte mit Open Source Software an und 
	ist dabei insbesondere auf Geodateninfrastrukturen, webbasierte Informationssysteme und WebMapping-Anwendungen fokussiert, 
	die neben 2D auch 3D Daten enthalten und auch für mobile Endgeräte optimiert sind.
	
	Wir orientieren uns dabei an den Anforderungen unserer Kunden und realisieren auf 
	den jeweiligen Bedarf zugeschnittene Lösungen.
	Neben der Anwendungsentwicklung können unsere Kunden Support und Schulungen für 
	verschiedene Open Source Softwaretools, z. B. PostgreSQL/PostGIS, UMN Mapserver, 
	GeoServer, SHOGun, OpenLayers, GeoExt, ExtJS und QGIS, erhalten.
	
	National und international gut vernetzt bieten wir zusammen mit unseren Partnerfirmen
	Wartung und Support zu Open Source Produkten/Projekten wie SHOGun, GeoServer oder OpenLayers an.
	Diese Rahmenbedingungen erlauben uns unvoreingenommen Lösungen zu entwickeln und 
	anzubieten, die tatsächlich alle Anforderungen unserer Kunden zufrieden stellen.
	}
}
\setlength{\fboxsep}{3pt}
\vfill

\newtimeslot{14:30}
% 14:30
% Der Untertitel ist ewig lang und der Inhalt steht auch so im Abstract.
\abstractEins{Sascha L. Teichmann}{MTSatellite}{Echtzeit Webmapping für Minetest-Welten oder Spiel-Spass mit GIS}{%
	Der Vortrag stellt die Freie Software MTSatellite vor, ein Live-Webmapping System für das Open World/Sandbox-Spiel Minetest 
	(eine Open-Source-Alternative zu Minecraft). 
	Der Vortrag führt das System vor und gibt eine teils vertiefende Übersicht über die eingesetzen Technologien sowohl aus GIS- 
	als auch aus Sicht eines passionierten Minetest-Spielers.}

\abstractZwei{Jakob Tworek}{Erfassung von Landnutzungsveränderungen mit FOSS Image Processing Tools}%
{%Change Detection mit QGIS - Landnutzungs- und -bedeckungsveränderungen erfassen
	}{%
	In diesem Vortrag wird eine Masterarbeit am Geographischen Institut der Universität Bonn vorgestellt. Im Rahmen dieser Masterarbeit 
	wurde zwei Change-Detection-Verfahren mit freier Bildverarbeitungssoftware angewandt und evaluiert.
	Zu den verwendeten Programmen zählen das Semi-Automatic Classification Plugin und Image Processing Tools der Orfeo Toolbox 
	und SAGA im Rahmen des Processing Framework von QGIS.
	}

\abstractZehn{Daniel Kastl}{Routing in der Datenbank}{Kürzeste-Wege-Berechnung und mehr mit pgRouting}{Die pgRouting-Erweiterung 
	ermöglicht es, auf Daten in einer PostgreSQL-Datenbank eine kürzeste-Wege-Suche und andere 
	netzorientierte Algorithmen anzuwenden. Neben den etablierten Funktionen sind für die nahe 
	Zukunft eine Reihe neuer Algorithmen zur Tourenplanung zu erwarten. Dieser Vortrag stellt 
	die verschiedenen Algorithmen vor und geht darauf ein, wie die Struktur der Netzdaten die Leistung des Systems beeinflussen kann.}

% Platz für Silber-Sponsor Camptocamp
\sponsorenbox{102_C2C-logo-baseline-RGB_PSE}{0.4\linewidth}{-1em}{%
	\textbf{Silbersponsor}\\
	Camptocamp ist ein innovatives Unternehmen im Bereich der Integration 
	von Software zum Aus\-werten von Geodaten, zur kompletten Verwaltung von 
	Unternehmen (ERP) und zur Verwaltung von Daten-Infrastrukturen.}

\newtimeslot{15:00}
% 15:00
%
\abstractEins{Tobias Knerr}{Straßenrennen in der Innenstadt}{SuperTuxKart-Strecken bauen mit OSM2World}{%
	Der Gründer des Open-Source-Projekts OSM2World erzählt von einer weiteren spannenden Anwendungsmöglichkeit des 3D-Renderers -- die 
	Erstellung von Rennstrecken für das freie Computerspiel SuperTuxKart.
	
	Der Vortrag beschreibt die Vorgehensweise, um die OSM-Daten einer Stadt in eine Rennstrecke zu 
	verwandeln. Voraussetzung ist eine ausreichend gute Abdeckung der jeweiligen Region mit 3D-Attributen in OpenStreetMap.
	
	Aus diesen Daten wird mit dem 3D-Renderer OSM2World ein realitätsnahes 3D-Modell erstellt, das optional 
	auch mit einem Geländemodell ergänzt werden kann. Mit dem 3D-Modelling-Tool Blender noch einige 
	Verbesserungen vorgenommen und der gewünschte Streckenverlauf eingezeichnet. Ein Blender-Plugin 
	erlaubt schließlich den Export einer fertigen Strecke für SuperTuxKart.	
}

\abstractZwei{Jens Schaefermeyer}{Automatisiertes Geodatenmanagement mit  GeoKettle}{}{%
	Dieser Vortrag stellt verschiedene Einsatzmöglichkeiten der freien ETL-Software GeoKettle vor. 
	GeoKettle ist die Open-Source-Alternative zur verbreiteten Software FME und kann nicht nur Geodatenformate konvertieren, 
	sondern beispielsweise auch Objekte verteilen und zusammenfassen, redundante Daten 
	finden oder Prozesse in einer grafischen Oberfläche modellieren.}

\abstractZehn{Felix Kunde}{PostGIS Memento}{Versionierung von PostGIS-Datenbanken}{Memento. Gibt es nicht 
	einen Film, der so heißt? Jemand, der kein Gedächtnis hat und sich alle Ereignisse aufschreibt? 
	Zumindest geht es bei pgMemento darum. pgMemento zeichnet alle Veränderungen in einer 
	PostgreSQL Datenbank auf und erlaubt die Wiederherstellung beliebiger frühere Zeitstände.}
%

\newtimeslot{16:00}
% 16:00
%
\abstractEins{Vladimir Elistratov}{Prokitektura -- prozedurale realistische 3D-Gebäude und Städte}{}{%
	Prokitektura is ein prozedurales und iteratives Verfahren für Schaffung architektonischer 3D-Modelle. 
	
	Ein Satz kleiner Funktionen auf Python wird verwendet um 3D-Gebäude zu generieren.
	Eine solche Funktion nennt man Regel. Jede nachfolgende Regel verfeinert das Modell 
	und ergänzt es mit zusätzlichen Details. Zur Zeit ist Prokitektura als ein Add-on für 
	Open-Source 3D-Plattform Blender realisiert.}

\abstractZwei{Till Adams}{WPS, GeoServer und SHOGun}{}{Der Vortrag stellt die Erweiterung des SHOgun 
	Frameworks als WPS-CLient dar. Als WPS Server kommt der GeoServer-WPS zum Einsatz. Der Vortrag 
	wird zum einen kurz die Mechanismen des WPS erläutern, die Einbindung in das SHOGun Framework an 
	praktischen Beispielen zeigen und am Ende die Möglichkeiten, die sich daraus für ein WebGIS ergeben vorstellen. }

\abstractZehn{Bjoern Schilberg}{Docker für den GIS-Einsatz}{}{Mit der leichtgewichtige Virtualisierungs-Software 
	Docker schnell und einfach Anwendungs-Container für Ihre GIS-Fachanwendung erstellen. Haben die 
	GIS-Admins demnächst wieder mehr Zeit für ungleich spannendere Aufgaben als die Installation von Fachanwendungen?}

% Platz für Silbersponsor Omniscale
% wir verzichten auf Absätze, da diese meist eh nur einen Satz lang sind 
% und die Sätze noch nicht ewig lang sind
\sponsorenbox{103_omniscale}{0.38\linewidth}{-1em}{%
	\textbf{Silbersponsor}\\
	Omniscale ist Spezialist für die Bereitstellung von schnellen und attraktiven Karten auf Ba\-sis von OpenStreetMap und amtlichen Daten.
	Unter maps.omniscale.com bietet Omniscale kostengünstige Kartendienste für 
	Desktop-GIS und Online-Anwendungen an. Zusätzlich entwickelt Omniscale für 
	Wirtschaft, Industrie und Behörden individuelle Kartendienste. Verlässlich, 
	schnell und attraktiv.
	Neben den Kartendiensten entwickelt Omniscale auch moderne Kartenanwendungen 
	auf Basis von Open-Source-Software. Omniscale ist zudem Initiator und Hauptentwickler 
	der Open-Source-Projekte MapProxy und Imposm.
}

\newtimeslot{16:30}
% 16:30
%
\abstractEins{Hongchao Fan}{Die neue WebGL-basierte Plattform für OSM-3D}{}{%
	In diesem Vortrag wird die technische Archtektur, das Funktiondesign und die 
	Implementierung der neue WebGL-basierte Plattform von OSM-3D präsentiert. 
	Es werden interaktive Visualisierungen der 3D-Stadtmodelle und 3D-Geländemodelle 
	demonstriert. Das neue Interface für interaktive Erfassungen von Dach- und 
	Fassadenstrukturen wird vorgestellt und mit Beispielen demonstriert. 
	Diskussion von Vorschlägen zur Verbesserungen der neuen Plattform. }

\abstractZwei{Astrid Emde}{Mapbender\,3 für den einfachen Aufbau von WebGIS Anwendungen}{So einfach ist 
	der Aufbau von WebGIS Anwendungen mit Mapbender3}{Mapbender3 ist eine Software zur einfachen 
	Erstellung von WebGIS Anwendungen. Der Vortrag geht vor allem auf die neuen Komponenten in 
	Mapbender3 ein und stellt diese vor. Anpassung an das eigne Layout, Suchen in eignen Daten und Rechteverwaltung werden thematisiert.}

% Der kurze Abstract ist auch gut.
%\abstractZehn{Volker Mische}{GeoCouch}{Ein multidimensionaler Index für Couchbase}{%
%	GeoCouch ist ein multidimensionaler Index für Couchbase, einer verteilten dokumentenorientierten Datenbank. 
%	Dieser Vortrag wird eine kurze Einführung in Couchbase, die räumliche Abfragen ermöglicht, 
%	und die Datenmodellierung geben und dann in eine Live-Demonstration übergehen. }
\abstractZehn{Volker Mische}{GeoCouch}{Ein multidimensionaler Index für Couchbase}{%
	Couchbase ist eine verteilte dokumentenorientierten Datenbank und gehört damit 
	zur Kategorie der NoSQL Datenbanken. Das bedeutet, dass das Datenmodell nicht 
	relational ist, also ein gewisses Umdenken bei der Strukturierung der Daten nötig ist. 
	GeoCouch ist in Couchbase integriert und ermöglicht räumliche Abfragen. Diese sind nicht 
	auf zwei Dimensionen begrenzt, sondern können durch weitere Attribute erweitert werden. 
	So ist es möglich Faktoren, wie Zeit, Ausmaß oder Kategorien einzubeziehen. Eine beispielhafte 
	multidimensionale Datenbankabfrage wären alle Einwohner eines bestimmten Gebiets, die unter 30 
	Jahre alt sind und in einer Mietwohnung mit weniger als 50m$^2$ wohnen. Dieser Vortrag wird eine 
	kurze Einführung in Couchbase und die Datenmodellerung geben und dann in eine 
	Live-Demonstration übergehen.}
% Platz für Bronzesponsor Büro Dassau
\sponsorenbox{209_gbd-consult}{0.34\linewidth}{-1.3em}{%
	\textbf{Bronzesponsor}\\
Die Geoinformatikbüro Dassau GmbH bietet Lösungen für die Umsetzung 
von GIS Projekten auf Open Source Basis und ist spezialisiert auf die 
Software GRASS und QGIS.
}

%

\newtimeslot{17:00}
% 17:00
\abstractEins{Christoph Hormann (Moderation)}{OSM Lightning Talks}{}{}\vspace{-2em}
\abstractZwei{Simon Jirka}{Erfahrungen mit Sensor Web-An\-wendungen für mobile Geräte und Desktop-Anwendungen}{}{%
	Dieser Vortrag stellt einen freien, auf JavaScript basierenden Sensor Web-Client vor, 
	der auf allen Arten von Endgeräten von Mobiltelefonen bis hin zu Desktop-PCs einsetzbar ist. 
	Er geht auf die Entwicklung des Clients und auf seine Nutzung als eigenständige Software und 
	als Basis eigener Projekte ein.}

\abstractZehn{Johannes Weskamm}{Vector-Tiles}{Performante Übertragung von umfangreichen Vektordaten}{%
	Es werden die Vor- und Nachteile von Vectortiling und der Vektordaten-Prozessierung im Client allgemein beleuchtet. 
	Als praktischer Anwendungsfall wird das Prinzip der Vectortiles anhand des Software-Stacks 
	TileStache, OpenLayers\,3 und PostGIS auf Basis von OSM-Daten vorgestellt.}
%

\newsmalltimeslot{17:30}
% 17:30
\vspace{-2em}
\abstractEins{Olaf Knopp}{PostNAS-Anwendertreffen}{}{%
	PostNAS ist ein Projekt zum Import von ALKIS-Daten über OGR. Es werden aktuelle Entwicklungen des Projekt vor\-gestellt.}
\abstractZwei{\label{bof-donnerstag}Till Adams}{SHOGun-BoF/Anwender- und Entwicklertreffen}{}{Es wird der aktuelle 
	Entwicklungsstand von SHOgun, neue Ideen und auch neue Entwicklungen, die seit der FOSSGIS 2014 dazu gekommen sind, vor\-gestellt.}
\abstractZehn{Danilo Bretschneider}{deegree-Anwendertreffen}{}{deegree ist ein Open-Source Java-Framework für 
	die Verwaltung und Darstellung OGC-konformer Geoinformationssysteme.}
%

\newsmalltimeslot{19:30}
% 19:00
%
%\abstractSenatssaal{}{Vollversammlung FOSSGIS e.V.}{}{}
	\thispagestyle{scrheadings} 
	\setlength\tabcolsep{0pt}
	\noindent\fcolorbox{white}{dezentrot}{\parbox{\titleboxwidth}{%
		\noindent\begin{tabu}{X[5L]r}
			& \talktime
			\tabularnewline
			\noindent\bfseries \normalsize \sectfont Vollversammlung FOSSGIS e.V.
			&
			Senatssaal
			\tabularnewline
		\end{tabu}%
	}}
	%
	\setpagebackground%
	\vspace{0.5em}% Abstand zum nächsten Talk, auch wenn es keinen Abstract gibt
	\setlength\tabcolsep{6pt} % Spaltenpadding wieder auf Default setzen
	
\sponsorenbox{207_intevation}{0.4\linewidth}{-1.3em}{%
	\textbf{Bronzesponsor}\\
Die Intevation GmbH ist ein produkt- und hersteller\-un\-ab\-hängiger IT-Dienstleister, 
welcher konsequent auf freie Software setzt. Mit unserer  bald fünfzehnjährigen praktischen 
Erfahrung haben wir eine Vielzahl von spannenden  Projekten mit den Themen Geographische 
Informationssysteme und  Webanwendungen, Netzwerk-Sicherheit, Email-Servern, 
Email-Kryptographie,  Groupware, Datenbanken, Sozialinformatik, Software-Paketierung  
und -Qualitätssicherung realisiert.
}
\vfill
\sponsorenbox{208_Kuestenschmiede}{0.34\linewidth}{-1.3em}{%
	\textbf{Bronzesponsor}\\
Wir entwickeln Webanwendungen, speziell zur Darstellung von Geoinformationen in 
Kombination mit Content Management Systemen (CMS). Dabei setzen wir auf verschiedene 
Open-Source-Projekte wie die OpenStreetMap,  OpenLayers, die Overpass API und auf das 
CMS Contao. Auf der FOSSGIS 2015 präsentieren wir als Aussteller unseren GIS-Baukasten 
con4gis, der es ermöglicht, auf einfache Weise komplexe, interaktive Geoinformationssysteme 
auf Webseiten bereitzustellen.
}

\newpage
\vfill
\noindent
\sponsorenbox{206_norBIT}{0.31\linewidth}{-1em}{%
	\textbf{Bronzesponsor}\\
Als Softwarehaus für GIS-Kom\-plett\-lösungen entwickelt die norBIT GmbH seit über 25 Jahren mit norGIS 
modular aufgebaute Fachschalen für Anwendungen im Bereich von Tiefbau, Ver- und Entsorgung, die mit 
verschiedenen Grafikplattformen, wie QGIS, AutoCAD und BricsCAD betrieben werden (ggf. auch gemischt).
Seit über 8 Jahren beteiligt sich die norBIT GmbH an der Weiterentwicklung des QGIS und nutzt die 
Software zur Erfüllung verschiedenster Kundenanforderungen von Auskunftslösungen bis zu GIS-Arbeitsplätze.
Mit dem norGIS-ALKIS-Import und dem ALKIS-Plugin für QGIS hat die norBIT GmbH eine Open-Source-Lösung 
für die Verarbeitung und Visualisierung von ALKIS-Daten in QGIS und UMN Mapserver geschaffen.
}
